% Priedų pavadinimas pridedamas prie pirmo priedo, be tuščio skiriamojo puslapio.
\clearpage
\phantomsection
\addcontentsline{toc}{section}{Priedai}
\priedas{Darbo planas ir jo pakeitimai}\label{app:planas}
\gaires{Tai siūloma darbo priemonė. Pildykite nuo projekto pradžios ir
išsaugokite ankstesnius susitarimus. Mokyklos nustatytos oficialios formos,
jei jos taikomos, pildomos atskirai.}
\begin{longtable}{@{}>{\raggedright\arraybackslash}p{2.9cm}>{\raggedright\arraybackslash}p{3.9cm}>{\raggedright\arraybackslash}p{3.9cm}>{\raggedright\arraybackslash}p{4.9cm}@{}}
\caption{Darbo planas}\label{tab:planas}\\
\toprule
Etapas & Planuotas terminas ir rezultatas & Faktinė data ir rezultatas & Pakeitimas, priežastis, aptarimas \\
\midrule\endfirsthead
\toprule
Etapas & Planuotas terminas ir rezultatas & Faktinė data ir rezultatas & Pakeitimas, priežastis, aptarimas \\
\midrule\endhead
Poreikio analizė & \pildyti{Data, rezultatas} & \pildyti{Data, įrodymas} & \pildyti{Įrašas} \\
Projektavimas & \pildyti{Data, rezultatas} & \pildyti{Data, įrodymas} & \pildyti{Įrašas} \\
Prototipo kūrimas & \pildyti{Data, rezultatas} & \pildyti{Data, įrodymas} & \pildyti{Įrašas} \\
Patikra ir tobulinimas & \pildyti{Data, rezultatas} & \pildyti{Data, įrodymas} & \pildyti{Įrašas} \\
Aprašo rengimas & \pildyti{Data, rezultatas} & \pildyti{Data, įrodymas} & \pildyti{Įrašas} \\
Pristatymo rengimas & \pildyti{Data, rezultatas} & \pildyti{Data, įrodymas} & \pildyti{Įrašas} \\
\bottomrule
\end{longtable}
\noindent\textbf{Su vadovu suderinti terminai}\par
Aprašo projektas autentiškumo patikrai: \pildyti{data}.\par
Galutinio produkto ir aprašo pateikimas: \pildyti{data}.\par
Pristatymas: \pildyti{data, vieta arba būdas}.

\priedas{Proceso ir konsultacijų žurnalas}\label{app:zurnalas}
\gaires{Kiekvienam reikšmingam etapui kopijuokite žemiau pateiktą įrašo formą.
Pridėkite nuorodą į tikrą įrodymą ir trumpą jo paaiškinimą. Vien darbų
saugyklos nuoroda nepakeičia proceso analizės.}
\noindent\textbf{Įrašas E1}\par
Data ir etapas: \pildyti{...}\par
Siektas rezultatas ir atliktas darbas: \pildyti{...}\par
Įrodymas (failas, versija, nuotrauka): \pildyti{...}\par
Klausimas vadovui ir gauta pastaba: \pildyti{...}\par
Priimtas sprendimas ir kitas žingsnis: \pildyti{...}\par
Trumpas įsivertinimas: \pildyti{kas pavyko, ką ir kaip tobulinti}.\par
\medskip
\noindent\textbf{Perimti komponentai ir pagalba}\par
\pildyti{Komponentas ar pagalbos priemonė; autorius ar teikėjas; versija;
nuoroda ir licencija; kur panaudota; ką pakeitėte ir kaip patikrinote.
Jei naudotas DI, pridėkite užklausos ar darbo sesijos nuorodą / išrašą, kai
to reikia pagal susitarimą su vadovu. Nerašykite, kad DI naudotas, jei jis nenaudotas.}

\priedas{Išsamūs patikros duomenys}\label{app:testai}
\gaires{Išsaugokite tikrus duomenis, taip pat nesėkmingus bandymus. Neskelbkite
dalyvių vardų ar kitų nereikalingų asmens duomenų. Jei testų daug,
pridėkite atskirą failą ir čia paaiškinkite jo struktūrą.}
\noindent\textbf{Bandymas T1}\par
Susijęs reikalavimas: \pildyti{R1}.\par
Data, produkto versija, aplinka: \pildyti{...}\par
Pradinė būsena, duomenys ir veiksmai: \pildyti{...}\par
Laukiamas rezultatas ir priėmimo kriterijus: \pildyti{...}\par
Gautas rezultatas ir įrodymo failas: \pildyti{...}\par
Vertinimas ir paaiškinimas: \pildyti{...}\par
Pataisa ir pakartotinis bandymas, jei reikėjo: \pildyti{...}.\par
\medskip
\pildyti{Pakartokite formą kitiems bandymams. Pridėkite matavimų lenteles,
apklausos klausimus ar kitą naudotą instrumentą, jei jis buvo taikytas.}

\priedas{Produkto pateikimas ir naudojimas}\label{app:produktas}
\noindent\textbf{Produkto identifikacija}\par
Pavadinimas ir galutinė versija: \pildyti{...}.\par
Pateikiami failai ir jų paskirtis: \pildyti{...}.\par
Saugykla / kopija / fizinio produkto vieta: \pildyti{...}.\par
\gaires{Skaitmeniniam produktui nurodykite versijos žymą arba pakeitimo ID,
priklausomybes ir paleidimo eigą. Fiziniam produktui -- surinkimo, veikimo ir
saugos sąlygas. Prieigą prie pateikto produkto patikrinkite iš kitos paskyros
ar įrenginio. Slaptažodžių ir API raktų apraše neturi būti.}
\noindent\textbf{Naudojimo žingsniai}
\begin{enumerate}
\item \pildyti{Reikalinga įranga ar programinė aplinka ir pasirengimas.}
\item \pildyti{Paleidimas arba surinkimas.}
\item \pildyti{Pagrindinės funkcijos naudojimas ir laukiamas rezultatas.}
\item \pildyti{Žinomi ribotumai, trikčių sprendimas ar saugus išjungimas.}
\end{enumerate}
\noindent\textbf{Demonstravimo planas}\par
\pildyti{Kokį trumpą scenarijų parodysite pristatydami produktą?
Kokią kopiją, vaizdo įrašą ar kitą atsarginį įrodymą turėsite sutrikus įrangai?}
